\documentclass[11pt,a4paper]{article}
\usepackage{curso}
\begin{document}
\portada{Lección}{2.1}{Movimiento browniano, proceso de Wiener y cálculo de Itô}

\noindent\textbf{Unidad 2 --- Dinámica estocástica.}\quad\textbf{Objetivos:} construir $W(t)$; entender $(dW)^2=dt$;
aplicar el lema de Itô. \textbf{Prerrequisitos:} 0.3; series de Taylor.

\section{Historia y motivación}
Brown (1827) observó el movimiento errático de granos de polen; Einstein (1905) lo explicó por choques moleculares y
predijo $\langle x^2\rangle=2Dt$ con $D=\kB T/(6\pi\eta a)$ (Einstein--Stokes). El movimiento browniano es el modelo
matemático del ruido térmico de toda la unidad \cite{gardiner}.

\section{El proceso de Wiener}
\begin{definicion}[ (proceso de Wiener)]
$W(t)$, $t\ge0$, con $W(0)=0$, trayectorias continuas, incrementos independientes y $W(t)-W(s)\sim\N(0,t-s)$ para $t>s$.
\end{definicion}
\textbf{Construcción:} una caminata con pasos de tiempo $\Delta t$ y saltos $\pm\sqrt{\Delta t}$ converge a $W$ cuando
$\Delta t\to0$ (teorema de Donsker). Si los saltos escalan como $\Delta t^\alpha$ con $\alpha>\frac12$ el límite es 0;
con $\alpha<\frac12$ explota. Propiedades: $\E[W(s)W(t)]=\min(s,t)$; autosemejanza $W(ct)\overset{d}{=}\sqrt c\,W(t)$;
continuo pero no derivable en ningún punto; el ruido blanco $\xi=\dot W$ solo existe como distribución.

\section{Variación cuadrática}
\begin{resultado}
Para particiones de $[0,T]$ con paso $\to0$, $\sum_k(\Delta W_k)^2\to T$ (en media cuadrática), mientras
$\sum_k|\Delta W_k|\to\infty$. Notación: $(dW)^2=dt$, $dW\,dt=0$, $(dt)^2=0$.
\end{resultado}
\begin{proof}[Idea]
$\E\sum(\Delta W_k)^2=\sum\Delta t_k=T$ y $\Var\sum(\Delta W_k)^2=2\sum\Delta t_k^2\to0$. En cambio
$\E\sum|\Delta W_k|=\sum\sqrt{2\Delta t_k/\pi}\to\infty$.
\end{proof}

\section{Integral y lema de Itô}
La integral de Itô $\int_0^tg(s)\,dW(s)$ se define con sumas evaluadas en el extremo izquierdo; para $g$ no anticipativa,
$\E\int g\,dW=0$ e $\E(\int g\,dW)^2=\int\E g^2\,ds$ (isometría de Itô).
\begin{resultado}[ (lema de Itô)]
Para $f$ dos veces derivable: $df(W)=f'(W)\,dW+\tfrac12f''(W)\,dt$. Para un proceso $dX=a\,dt+b\,dW$:
$df(X)=\big(f'a+\tfrac12f''b^2\big)dt+f'b\,dW$.
\end{resultado}
\begin{ejemplo}
$d(W^2)=2W\,dW+dt$ $\Rightarrow$ $\E[W(t)^2]=t$. $d(e^W)=e^W\,dW+\frac12e^W\,dt$ $\Rightarrow$ $\E[e^{W(t)}]=e^{t/2}$.
\end{ejemplo}
\textbf{Itô frente a Stratonovich.} Con ruido aditivo (coeficiente constante) coinciden; con ruido multiplicativo
difieren por un término de deriva $\frac12bb'$. En este curso usamos Itô.

\section{Ejemplo físico}
Esfera de $a=1\ \mu$m en agua ($\eta=10^{-3}$ Pa\,s) a 300 K: $D=4{,}14\times10^{-21}/(6\pi\cdot10^{-3}\cdot10^{-6})\approx2{,}2\times10^{-13}$ m$^2$/s;
$x(t)=\sqrt{2D}\,W(t)$; $x_{\text{rms}}(1\ \text{s})\approx0{,}66\ \mu$m; en 100 s, solo 10 veces más.

\begin{aplicacion}
Fuente de ruido de las ecuaciones de Langevin (2.2), de los modelos de difusión (Unidad 9) y del solucionador OU (L3).
\end{aplicacion}

\begin{error}
\begin{itemize}[nosep]
\item Saltos $\propto\Delta t$ en vez de $\sqrt{\Delta t}$.
\item Regla de la cadena ordinaria (falta $\frac12f''dt$).
\item Tratar $\xi(t)$ como función ordinaria.
\item Mezclar Itô y Stratonovich con ruido multiplicativo.
\end{itemize}
\end{error}

\section{Resumen}
$W$: incrementos $\N(0,t-s)$; $(dW)^2=dt$; Itô: $df=f'dW+\frac12f''dt$; difusión $\propto\sqrt t$.

\begin{thebibliography}{9}\small
\bibitem[1]{gardiner} C. W. Gardiner, \emph{Handbook of Stochastic Methods}, Springer.
\bibitem[5]{kloeden} P. E. Kloeden y E. Platen, \emph{Numerical Solution of Stochastic Differential Equations}, Springer, 1992.
\bibitem[46]{holderrieth} P. Holderrieth y E. Erives, An introduction to flow matching and diffusion models, 2025.
\end{thebibliography}
\end{document}
